\section{从磁共振信号到临床图像：技术演进的共同基础}

\subsection{射频激发、弛豫与组织对比}

MRI的起点不是直接给人体拍照，而是测量原子核在磁场中的响应。人体内氢质子数量多，在主磁场$B_0$中形成可观测的宏观磁化，并按拉莫尔频率进动。忽略方向时，两者的关系可写为
\begin{equation}
    \omega_0=\gamma B_0,
    \label{eq:larmor}
\end{equation}
其中，$\gamma$为旋磁比。施加与拉莫尔频率一致的射频脉冲，磁化矢量偏离平衡方向，横向分量出现，接收线圈这才有信号可收。脉冲一撤，系统通过$T_1$纵向弛豫把纵向磁化拉回平衡，通过$T_2$横向弛豫丢掉横向相干性，自由感应衰减信号就是这么来的\cite{wu2024mri}。

组织之间质子密度和弛豫时间不同，调整重复时间、回波时间这类序列参数，就能把某一种对比突出来。T1WI和T2WI最先回答的仍是位置、边界以及与周围结构的关系。这里已经能看出MRI与普通形态摄影的区别：灰度不是组织密度的直接翻版，它是物理状态与序列设计共同作用的结果。初学时容易略过这一点，而后面DWI、DCE-MRI、fMRI名称各异，底下做的其实是同一件事，即对磁共振信号作有选择的编码。

\subsection{梯度空间编码与K空间重建}

只有自由感应衰减信号，得到的是整段组织叠在一起的总响应，判断不了信号来自何处。梯度磁场让不同位置的局部磁场和进动频率产生可控差异，层面选择、相位编码与频率编码由此给信号安上空间坐标。时域信号按特定轨迹填入K空间，再经傅里叶变换重建为图像。教材把这件事讲得很直接：成像信号方程揭示的底层逻辑，就是用采集到的时域信号反演空间体素的磁化属性\cite{wu2024mri}。所以K空间不是一张“中间图片”，它是空间频率数据的排布方式。采样范围管视场，采样间隔管分辨率，轨迹则同时牵动扫描时间和伪影。

\begin{figure}[ht]
    \centering
    \includegraphics[width=.92\linewidth]{textbook/mri-signal-kspace}
    \caption{磁共振信号采集、空间编码与图像重建的基本流程\cite{wu2024mri}}
    \label{fig:mri-signal}
\end{figure}

从工程角度看，梯度系统不只管空间定位。换个施加方式，图像就能对扩散、血流这些过程敏感；快速读出技术成熟之后，动态过程也进得来。平面回波成像提高了采集速度，DWI和BOLD-fMRI才有了现实条件\cite{wu2024mri,sun2001fmri}。MRI的技术演进因此不是另起炉灶，它是在“激发—编码—采集—重建”这条既有链条上，一处一处加进新的信息维度。

\subsection{从形态对比到多参数信息：本文的技术主线}

把前面的内容串起来，本文把技术演进分成三层。底下的第一层是形态显示，T1WI、T2WI把病灶的位置和结构关系说清楚；往上的第二层是功能与定量表征，DWI/ADC、DTI/DKI、DCE-MRI和BOLD-fMRI把扩散、微循环及脑活动变成可分析的数字；再往上才轮到决策支持，多个参数与临床特征、影像组学、机器学习结合，用于病理预测、疗效评价与风险评估。上面两层都离不开底下这一层，功能参数再漂亮，也要有准确的解剖定位和过得去的图像质量兜底。

\begin{figure}[ht]
    \centering
    \includegraphics[width=.90\linewidth]{selfdrawn/mri-tech-tree.png}
    \caption{磁共振由形态显示走向定量评估与决策支持的技术演进树}
    \label{fig:tech-tree}
\end{figure}

\section{扩散磁共振：从水分子运动到组织微结构}

\subsection{DWI与ADC：把扩散过程转化为图像对比}

水分子在组织里从不静止，随机热运动一直在进行。DWI在序列中加入扩散敏感梯度，让发生位移的自旋产生相位分散，运动越明显，信号衰减通常越厉害。$b$值描述加权的程度：$b$值越大，对扩散越敏感，信噪比也越低。用不同$b$值下的信号，可以估出表观扩散系数（apparent diffusion coefficient，ADC）。在简化的单指数模型下，其关系可表示为
\begin{equation}
    S(b)=S_0\exp(-b\cdot \mathrm{ADC}),
    \label{eq:adc}
\end{equation}
其中，$S_0$表示不考虑扩散衰减时的信号\cite{wu2024mri,zeng2005dwi}。所以DWI图和ADC图给出的并不是一回事：前者是加权图像，后者是在模型假设下算出来的定量参数。只看“DWI高信号”就断定扩散受限，容易出错。

急性脑梗死是扩散成像最典型的用武之地。高燕等在小体积急性梗死的研究中以DWI结果为参照，比较CTA、CTP及其联合方案，联合方案的诊断效能高于各自单独使用，这也从侧面说明DWI在认出梗死灶这件事上的分量\cite{gao2025infarct}。不过重点并不在“DWI可以替代所有检查”，而在三种技术各自回答不同的问题：梗死灶在哪、血管窄不窄、灌注够不够。扩散信息进入临床流程，改变的是疾病分层的思路，其他检查的必要性并没有因此取消。

\subsection{从DWI到DTI、DKI：模型如何逐步接近组织复杂性}

单一ADC把组织扩散压成一个标量，可真实组织既有方向，也有复杂的微环境。扩散张量成像（diffusion tensor imaging，DTI）用多个方向的扩散编码描述各向异性，能估出主要扩散方向，并进一步得到各向异性分数（FA）、平均扩散率（MD）等指标和纤维束追踪结果。问题于是从“扩散快不快”变成“往哪个方向扩、被什么结构挡着”。DTI目前已被用于创伤性颅脑损伤、中风、脑肿瘤及神经退行性疾病等研究，其指标仍受采集与分析条件影响\cite{zeng2005dwi,wang2024dti}。

DWI和常规DTI都默认扩散是高斯的，而细胞膜、细胞器及多隔室环境会让真实扩散偏离这个理想模型。扩散峰度成像（diffusion kurtosis imaging，DKI）在DTI基础上描述非高斯扩散，常用平均峰度（MK）表征微结构的复杂程度\cite{zeng2011dki}。这一步演进说明，技术进步常常表现为模型一点一点逼近现实，跟图像外观是否“高级”关系不大。代价是模型越复杂，需要的采集方向越多，对信号质量和拟合条件也越挑剔。DKI并非在所有场合都优于DTI或DWI。

\subsection{案例：由扩散参数推测肿瘤侵袭性}

周定燕等回顾性分析了66例前列腺癌患者，把T2WI、DWI/ADC与DCE-MRI参数合起来，预测Ki-67表达和Gleason评分。ADC与两者均呈负相关，联合模型预测Ki-67表达的AUC达0.940，预测Gleason评分为0.861\cite{zhou2024prostate}。水分子扩散与肿瘤增殖、分级就这样被联系起来，扩散成像也从检出病灶延伸到病理表型评估。

只是这种联系得当作有病理依据的统计关联来看，而不是ADC在直接观察细胞行为。ADC偏低，可能来自细胞密度，也可能来自间质结构或其他因素，不能反推所有低ADC病灶都有同样的病理性质。对我来说，分清模型参数、组织微结构和临床病理结论这三层，是防止过度解释的第一步。

\subsection{技术限制与改进方向}

扩散成像要强梯度，也要快采集，麻烦集中在运动伪影、涡流、EPI几何畸变和低信噪比上。$b$值取多少、方向数取多少、用什么拟合方法，都会动到定量结果\cite{wu2024mri}。所以这条路往前走，不能只顾把微结构模型做得更复杂，硬件、校正算法和采集规范得跟上。ADC的好处恰恰在于模型简单、临床做得起来；DTI和DKI则应当用在确有增量信息的地方。让复杂度与问题相匹配，比单纯比参数个数实在。

\section{DCE-MRI：从强化形态到微循环定量}

\subsection{动态采集与药代动力学分析}

常规增强MRI看的是某一时刻病灶有没有强化、强化成什么形态。DCE-MRI在对比剂注入前后连续采集，把组织信号随时间的变化记下来。时间—信号强度曲线反映对比剂到达、聚集和消退的过程，可以定性看，也可以作半定量或定量分析。王艾博和边杰指出，DCE-MRI突破了常规MRI主要分析病变形态的限制，配合定量、半定量分析与药代动力学模型，能从微循环层面给出更完整的信息\cite{wang2016dce}。

常用参数里，容积转运常数$K^{\mathrm{trans}}$描述对比剂由血浆向血管外细胞外间隙转运的综合过程，速率常数$K_{\mathrm{ep}}$描述回流，$V_{\mathrm{e}}$是血管外细胞外间隙的容积分数。有了它们，讨论就能从“亮不亮”推进到血流供应、血管通透性及组织间隙特征。但这些参数依赖具体的生理模型和测量条件，不是脱离模型的纯生物学常数\cite{wu2024mri,zhou2024prostate}。

\subsection{案例一：乳腺肿瘤诊断中的价值与证据边界}

马国津等把DCE-MRI、超声和钼靶X线放在一起比较，纳入266例乳腺肿瘤患者，其中47例做了DCE-MRI。报告给出的灵敏度为100.00\%、特异度为85.71\%，与病理诊断一致性较高；同时钼靶对微小钙化灶的检出率更高，超声联合钼靶的诊断效能也不错\cite{ma2024dce}。DCE-MRI确实提供了别的检查给不了的血供信息，但据此说它该替代其他影像，就走远了。还有个细节不能忽略：做DCE-MRI的47例只是总样本的一部分，高性能要站得住，得在更大、检查人群更统一的研究里再看。

另一项150例乳腺癌研究里，腋窝淋巴结转移组的MK、$K^{\mathrm{trans}}$和$K_{\mathrm{ep}}$偏高，MD偏低；DCE-MRI模型AUC为0.833，高于DKI模型的0.733，联合模型在灵敏度和准确度上又改善了一些\cite{an2024dki}。这说明“联合”的价值不在把序列机械相加，而在于扩散与微循环信息可能互补。反过来也提醒一句：某一参数在特定任务上表现突出，并不意味着更复杂的模型在各项指标上都占优。

\subsection{从参数到决策仍需跨越的环节}

DCE-MRI的定量能力建立在动态采集、运动控制、信号—浓度转换及模型拟合之上。时间分辨率不够，曲线细节会漏；空间分辨率不够，小病灶和感兴趣区就勾不准；分析模型与扫描方案一换，参数可比性又成问题。对比剂、扫描时长、患者配合程度，都得算进临床实施成本。

我读这些材料时最有感触的一点是，定量并不等于自动客观。参数以数字出现，确实比肉眼描述少些模糊，可数字背后仍是采集条件和模型假设。要让$K^{\mathrm{trans}}$这类参数进入治疗决策，光证明组间有差异不够，还得说清阈值能不能复现、在常规影像之外是否真的多给了有用信息、多出来的扫描与分析成本值不值。

\section{功能磁共振：从血氧变化到脑网络}

\subsection{BOLD效应与快速成像}

fMRI记录的不是神经元放电。它借的是神经活动、能量代谢与局部血流之间的那层关系：神经活动增强，局部脑血流增加幅度通常超过氧消耗增加幅度，脱氧血红蛋白因此相对减少。脱氧血红蛋白有顺磁性，其变化会影响局部磁场均匀性和$T_2^*$加权信号，血氧水平依赖（blood oxygenation level-dependent，BOLD）对比由此形成\cite{wu2024mri,sun2001fmri}。说到底，BOLD是经神经血管耦合转了一道手的间接信号。

BOLD信号弱，还随时间变化，需要快速、连续的全脑采集。EPI这类快速成像技术让刺激与信号变化能在时间上对上，fMRI才走得动。孙学军等的早期综述提到，fMRI受快速成像技术影响，由有创走向无创，EPI与快速小角度激发技术在时间和空间分辨率上各有长短\cite{sun2001fmri}。可见fMRI的突破一半来自生理学认识，一半来自采集速度这个工程条件。

\subsection{从任务态定位到静息态脑网络}

任务态fMRI设计特定刺激或行为任务，去找与任务相关的脑区；静息态fMRI不设明确任务，看的是自发脑活动。后者衍生出低频振幅（ALFF）、局部一致性（ReHo）和功能连接等指标：ALFF看特定频段的信号强度，ReHo看局部体素时间活动是否一致，功能连接看不同脑区时间信号的相关关系\cite{wu2024mri}。问题也就从“哪个脑区被激活”变成“多个脑区怎么协同”。

一项首发重度抑郁症研究纳入40例患者与50名健康对照，静息态fMRI显示感觉运动、视听觉、默认网络和执行控制网络存在广泛异常，作者把这些结果定位为理解疾病发生发展的初步线索\cite{ye2024mdd}。阿尔茨海默病综述则认为，rs-fMRI为寻找早期无创标记物提供了可能\cite{he2024ad}。这些研究让fMRI显得有吸引力，但“可能提供线索”与“已经能够独立诊断”之间的距离，并不小。

\subsection{案例：康复疗效与脑功能重塑能否相互证明}

罗红和徐丽把40例脑出血患者分成中枢、外周和联合三组，在常规治疗基础上分别给予重复经颅磁刺激、重复外周磁刺激或二者联合。联合组上肢功能改善更明显，多个脑区与运动皮层的功能连接也发生变化，部分连接强度的变化与FMA-UE评分变化相关\cite{luo2024rehab}。这项随机对照研究把临床量表和影像指标放进同一框架，比单讲治疗前后脑图差异要有说服力。

不过作者解释机制时用了“可能由于”，这个措辞值得保留。量表改善与功能连接相关，能支持脑功能重塑的假说，却不足以证明某一处连接变化就是功能恢复的直接原因。fMRI在疗效研究里比较合适的位置，是提供机制线索和候选观察指标；要确证中介机制，纵向设计、对照和重复验证一样都不能少。

\subsection{功能影像的证据限度}

fMRI容易受头动、呼吸和心血管生理波动影响，预处理流程与统计阈值也能改变结果。个体之间解剖和功能的差异，又使群体层面的显著差异很难直接搬到个体诊断上。针灸影像学研究中关于样本量的讨论有普遍意义：样本量估算随意，结果的信效度就受损，研究者得按具体fMRI方案保证统计效力，并推动数据共享\cite{zhou2024sample}。

比来比去，我的判断反而更谨慎了。脑网络图看着直观，可它的解释链条比常规形态图像长得多：神经活动、血流变化、MR信号、预处理、统计模型、临床行为，每一环都可能加进不确定性。fMRI最打动人的地方，是它能提出过去问不出的问题，而不在于它已经替所有脑疾病备好了答案。

\section{从多参数成像到临床决策支持}

\subsection{多参数的价值在于互补，而不只是数量}

前面几条线索分别描述形态、扩散、微循环与脑功能。多参数MRI的意义，在于借这些信息的互补性去回答单一序列答不了的问题。前列腺癌那个案例里，DCE-MRI药代动力学参数与ADC联合后，对Ki-67和Gleason评分的预测优于单项指标\cite{zhou2024prostate}。王雨薇等基于246例患者建立DCE-MRI与DWI的瘤内、瘤周联合模型，训练组和验证组AUC分别是0.953和0.948，再加上肿瘤最大径后效能更高\cite{wang2024her2}。另一项401例研究的多参数模型，训练集与测试集AUC为0.932和0.906\cite{liu2024multiparam}。多源信息确实有机会把影像从诊断工具推向病理表型和风险预测。

但联合模型不必然赢过简单模型。张丁懿等比较DWI、DCE及二者联合的乳腺癌分子分型模型，测试组AUC依次为0.722、0.798和0.821；除训练组DCE与联合模型的比较外，其余两两差异均无统计学意义\cite{zhang2024luminal}。评价多参数技术时只报联合模型的最高值没有意义，得在相同数据和参照标准下与简单模型比，再说清这点提升稳不稳、值不值那份额外成本。

\subsection{影像组学与人工智能属于决策层，而非新的信号来源}

影像组学做的事，是从既有图像中提取形状、强度、纹理这类高维特征，再经筛选和建模去预测病理或结局；机器学习与深度学习提供了更多特征学习和模式识别的手段。它们改变的是图像分析方式，氢质子的物理响应一点没变。所以本文把它们放在“支持决策”这一层，不与DWI、DCE-MRI并列成新的采集技术。这样分层的好处是误差来源比较清楚：采集阶段带进来的噪声和伪影，后续算法不会自动清掉，复杂模型甚至可能把这些差异学成疾病特征。

失｜组织成分与血流动力学｜脂肪肝、颈动脉狭窄｜定量条件及验证｜无创定量与风险评估。

\begin{table}[ht]
    \centering
     \includegraphics[width=.90\linewidth]{selfdrawn/mri-technique-comparison}
    \caption{主要磁共振技术及图像分析方法的功能与局限比较}
    \label{tab:tech-compare}
\end{table}

\subsection{定量技术的拓展：从“参数更多”走向“问题更明确”}

除三条主线之外，MRI-PDFF和4D Flow MRI展示了定量成像的另一面。MRI-PDFF把肝脂肪从视觉分级推进到无创定量，相关综述强调了它的精确性和可重复性\cite{yang2024pdff}。4D Flow MRI用血流动力学参数看颈动脉狭窄，一项20例患者与26名对照的研究发现，能量损失有助于诊断狭窄，低壁剪切力及狭窄下游压力梯度可能与急性缺血性卒中相关\cite{lu2024flow}。定量MRI不只有肿瘤一个出口，组织成分测量和血流风险评估同样用得上。

这些例子又一次说明，有临床意义的技术未必是最复杂的技术。一个参数若能稳定回答一个明确问题，可能比塞满特征却难以解释的模型更有价值。技术演进既包含向更复杂的模型前进，也包含把某一次测量做得更标准、更可靠、更容易实施。

\subsection{生物医学工程专业学生视角下的思考}

最后我想谈谈自己的三点思考。

最要紧的一点，是物理上的可测量，不等于临床上的可解释。DWI测的是扩散相关的信号衰减，DCE-MRI测的是对比剂的动态分布，BOLD测的是血氧相关的信号变化，三者都要过一道模型，才能和细胞密度、血管通透性或神经活动接上头。工程研究如果只盯着最终图像好不好看，中间假设是否成立就没人追问了。这也是我觉得本专业学生得同时学物理、统计和医学的原因：缺哪一环，都可能把“测得到”当成“已经知道”。

第二点是老生常谈，却常常被跳过：指标更定量，结果不一定更可重复。定量参数依赖设备、序列、校正、感兴趣区和分析方法。脑小血管病MRI专家共识之所以强调扫描和报告规范，是为了让不同实践场景下的结果可以相互比较\cite{csvd2024consensus}。规范化在我看来不是给创新设限，它是让创新成果能被别人检验的前提。新序列只在单一设备和单一团队里有效，临床价值终究有限。

第三点关于代价。新增序列要占扫描时间，新增算法要背上验证负担。合理的研究设计应当先锁定临床问题，用简单方法立好基线，再证明复杂技术究竟多带来了什么。本文引用的研究多数是单中心、回顾性或探索性的，有些AUC看着漂亮，仍需独立人群和前瞻性研究来检验。我对MRI研究方向的兴趣，大概就来自这条还没闭合的证据链：工程师不能只问“能不能做到”，还得和临床一起回答“可不可靠、要不要用、患者是否真的受益”。

写到这里，想再说几句课程本身的事。这门课给我的收获比我预想的大。大一的时候我对学院的印象基本停在“生医工好像什么都做”这一层，具体有哪些方向、它们之间怎么区分，其实说不清。到这门课上，磁共振成像这条线被一点点拆开：从拉莫尔进动到K空间，从扩散加权到脑网络，一个研究方向第一次在我眼前显出完整轮廓。选磁共振做论文题目，算是顺着自己的手感走的，因为上课前我就对这个研究方向有兴趣，但那种兴趣更接近“觉得挺神”，说不清神在哪儿。写的过程里读了些近几年的文献，把DWI、DCE、fMRI三条线各自从哪儿长出来、现在卡在什么地方捋了一遍，兴趣反倒踏实了。那些好看的图像背后，其实是模型假设、采集条件、验证设计这些琐碎却决定成败的东西。所以这门课我愿意推荐给别人，尤其是大一的学弟学妹。晚一年听当然也不迟，但早点接触，至少能早点摸清自己感兴趣的研究方向。我自己的体会是，兴趣得先落到具体题目上，方向感才长得出来，规划也才有意义。
